\chapter{HCG Levels}
\index{HCG Levels}

\section*{Definition and Overview}
HCG (human chorionic gonadotropin) is a hormone produced by the placenta shortly after a fertilised egg implants in the uterus. It is the hormone detected by pregnancy tests, and its level in the blood rises rapidly in early pregnancy. HCG levels are measured to confirm pregnancy, to monitor the progress of early pregnancy, and to help diagnose conditions such as ectopic pregnancy, miscarriage, and gestational trophoblastic disease. The expected levels of HCG vary widely between women and pregnancies, and a single measurement is less informative than the trend over time. In a normal early pregnancy, HCG roughly doubles every 48--72 hours.

\section*{Epidemiology}
HCG testing is one of the most common laboratory tests in many countries. It is used in the diagnosis of virtually all pregnancies, in the management of early pregnancy problems (which affect around 1 in 5 confirmed pregnancies), and in the diagnosis and follow-up of ectopic pregnancies, which occur in about 1 in 80--100 pregnancies. HCG measurements also guide the management of molar pregnancy and are used as a tumour marker in some cancers. Home pregnancy tests detect HCG in urine and are used by millions of people.

\section*{Aetiology and Pathophysiology}
HCG levels reflect the activity of the developing placenta.

\begin{itemize}
    \item \textbf{Production:} HCG is produced by the syncytiotrophoblast cells of the early placenta from the time of implantation
    \item \textbf{Timing:} HCG becomes detectable in the blood around 8--11 days after conception and in urine shortly after
    \item \textbf{Rise in normal pregnancy:} HCG doubles approximately every 48--72 hours in the first weeks, peaking around 8--10 weeks, then falling and stabilising
    \item \textbf{Wide normal range:} normal HCG levels vary widely between women and are influenced by the number of fetuses
    \item \textbf{Ectopic pregnancy:} HCG rises more slowly and irregularly when the pregnancy is outside the uterus
    \item \textbf{Miscarriage:} falling or plateauing HCG suggests a failing pregnancy
    \item \textbf{Molar pregnancy:} HCG levels are often very high
    \item \textbf{Tumour marker:} some cancers produce HCG
\end{itemize}

\section*{Clinical Features}
HCG levels are measured in specific clinical situations.

\begin{itemize}
    \item Confirming and dating an early pregnancy
    \item Investigating bleeding or pain in early pregnancy
    \item Suspected ectopic pregnancy, where serial levels and ultrasound are used
    \item Suspected miscarriage
    \item Monitoring after treatment of ectopic or molar pregnancy
    \item Screening for Down syndrome and other conditions as part of combined screening
    \item Follow-up of certain cancers
\end{itemize}

\section*{Differential Diagnosis}
The pattern of HCG results must be interpreted alongside other information.

\begin{itemize}
    \item Normal intrauterine pregnancy
    \item Ectopic pregnancy
    \item Miscarriage or failing pregnancy
    \item Molar (hydatidiform) pregnancy
    \item Multiple pregnancy, which produces higher levels
    \item Incorrect dating of the pregnancy
    \item HCG-producing tumours in non-pregnant people
    \item Home test errors and urine dilution
\end{itemize}

\section*{When to Seek Medical Care}
See a doctor for a positive pregnancy test for antenatal care, and promptly for bleeding, pain, or other concerns in early pregnancy.

\textbf{Contact your local emergency services for:}
\begin{itemize}
    \item Severe abdominal or pelvic pain with a positive pregnancy test, which may indicate ectopic pregnancy
    \item Heavy vaginal bleeding, especially with faintness or dizziness
    \item Shoulder-tip pain with bleeding and a positive pregnancy test
    \item Signs of shock, including pallor, rapid pulse, and collapse
\end{itemize}

\section*{Diagnosis and Investigations}
HCG is measured in blood or urine, and results are interpreted with ultrasound.

\begin{itemize}
    \item \textbf{Urine pregnancy test:} detects HCG for home and clinical pregnancy testing
    \item \textbf{Blood HCG:} quantitative measurement of the exact level
    \item \textbf{Serial measurements:} two or more levels 48--72 hours apart show the trend
    \item \textbf{Ultrasound:} visualises the pregnancy and distinguishes intrauterine from ectopic
    \item \textbf{Interpretation:} results are interpreted with the gestational age, symptoms, and ultrasound findings
    \item \textbf{Follow-up:} repeated measurements until the situation is resolved
\end{itemize}

\section*{Management}
Management depends on what the HCG pattern indicates.

\begin{itemize}
    \item \textbf{Normal pregnancy:} reassurance and routine antenatal care
    \item \textbf{Uncertain early pregnancy:} serial HCG and repeat ultrasound until the diagnosis is clear
    \item \textbf{Ectopic pregnancy:} treatment with medicine (methotrexate) or surgery, with HCG monitoring until it returns to zero
    \item \textbf{Miscarriage:} supportive care, and monitoring until HCG falls
    \item \textbf{Molar pregnancy:} surgical evacuation and regular HCG monitoring
    \item \textbf{Monitoring after treatment:} HCG is tracked to ensure complete resolution
    \item \textbf{Counselling:} emotional support for the outcomes
\end{itemize}

\section*{Complications}
\begin{itemize}
    \item Undiagnosed ectopic pregnancy, which can rupture and cause life-threatening bleeding
    \item Missed miscarriage with retained tissue
    \item Gestational trophoblastic disease, which can persist and needs monitoring
    \item Anxiety and distress from uncertain results
    \item Misinterpretation of a single HCG level
    \item Complications of treatment for ectopic or molar pregnancy
\end{itemize}

\section*{Prognosis and Outcomes}
For a normal pregnancy, HCG rises predictably, and outcomes are excellent. When HCG patterns are abnormal, they guide the diagnosis and management of ectopic pregnancy, miscarriage, and molar pregnancy, allowing timely treatment that preserves health and fertility. Most ectopic pregnancies are now diagnosed early through HCG and ultrasound, reducing mortality. After treatment, HCG returns to zero, confirming resolution. With appropriate care, most women go on to have healthy future pregnancies.

\section*{Prevention and Self-Care}
\begin{itemize}
    \item Attend antenatal care early in pregnancy
    \item Follow the schedule of blood tests and scans recommended by your maternity team
    \item Seek prompt medical advice for bleeding or pain in early pregnancy
    \item Attend all follow-up HCG tests after an ectopic or molar pregnancy
    \item Use a reliable test and confirm positive results with a doctor
    \item Discuss family planning after treatment of ectopic or molar pregnancy
    \item Seek emotional support after miscarriage or ectopic pregnancy
\end{itemize}

\section*{Sources and Further Reading}
The following reputable sources provide further information for healthcare students and patients seeking reliable, current advice about HCG levels.

\begin{itemize}
    \item Royal Australian and New Zealand College of Obstetricians and Gynaecologists. \textit{Early pregnancy information.} https://ranzcog.edu.au/
    \item Pregnancy, Birth and Baby (Australian Government). \textit{Pregnancy tests and HCG.} https://www.pregnancybirthbaby.org.au/
    \item Healthdirect Australia. \textit{HCG blood test.} https://www.healthdirect.gov.au/
    \item Royal College of Pathologists of Australasia. \textit{Patient information on HCG.} https://www.rcpa.edu.au/
    \item NHS. \textit{Human chorionic gonadotrophin (HCG) test.} https://www.nhs.uk/
    \item MedlinePlus. \textit{HCG blood test.} https://medlineplus.gov/
\end{itemize}
